% Header aus der Vorlage
\documentclass[a4paper,11pt, footheight=26pt]{scrartcl}
\usepackage[head=23pt]{geometry}	% head=23pt umgeht Fehlerwarnung, dafür größeres "top" in geometry
\geometry{a4paper, top=30mm, left=20mm, right=20mm, bottom=22mm,headsep=10mm, footskip=12mm}

% Input inkl. Umlaute, Silbentrennung
\usepackage[T1]{fontenc}
\usepackage[utf8]{inputenc}
\usepackage[ngerman]{babel}
\usepackage{csquotes}	% Anführungszeichen

% HTW Corporate Design: Arial (Helvetica)
\usepackage{helvet}
\renewcommand{\familydefault}{\sfdefault}

% Style-Aufhübschung
\usepackage{soul, color}	% Kapitälchen, Unterstrichen, Durchgestrichen usw. im Text
\usepackage{scrlayer-scrpage}	% Kopf-/Fußzeile
%\usepackage{titleref}
\usepackage[perpage]{footmisc}	% Fußnotenzählung Seitenweit, nicht Dokumentenweit

% Mathe usw.
\usepackage{amssymb}
\usepackage[fleqn]{amsmath}	% fleqn: align-Umgebung rechtsbündig
\usepackage{xcolor}
\usepackage{esint}	% Schönere Integrale, \oiint vorhanden
\everymath=\expandafter{\the\everymath\displaystyle}	% Mathe Inhalte werden weniger verkleinert
\usepackage{wasysym}	% mehr Symbole, bspw \lightning

% tikz usw.
\usepackage{tikz}
\usepackage{pgfplots}
\pgfplotsset{compat=1.11}	% Umgeht Fehlermeldung
\usetikzlibrary{graphs}
%\usetikzlibrary{through}	% ???
\usetikzlibrary{arrows}
\usetikzlibrary{arrows.meta}	% Pfeile verändern / vergrößern: \draw[-{>[scale=1.5]}] (-3,5) -> (-3,3);
\usetikzlibrary{automata,positioning} % Zeilenumbruch im Node node[align=center] {Text\\nächste Zeile}
\usetikzlibrary{patterns}	% Schraffierte Füllung
\tikzstyle{reverseclip}=[insert path={	% Inverser Clip \clip
	(current page.north east) --
	(current page.south east) --
	(current page.south west) --
	(current page.north west) --
	(current page.north east)}
% Nutzen: 
%\begin{tikzpicture}[remember picture]
%\begin{scope}
%\begin{pgfinterruptboundingbox}
%\draw [clip] DIE FLÄCHE, IN DER OBJEKT NICHT ERSCHEINEN SOLL [reverseclip];
%\end{pgfinterruptboundingbox}
%\draw DAS OBJEKT;
%\end{scope}
%\end{tikzpicture}
]	% Achtung: dafür muss doppelt kompliert werden!
\usepackage{graphpap}	% Grid für Graphen

% Tabular
\usepackage{longtable}	% Große Tabellen über mehrere Seiten
\usepackage{multirow}	% Multirow/-column: \multirow{2[Anzahl der Zeilen]}{*[Format]}{Test[Inhalt]} oder \multicolumn{7[Anzahl der Reihen]}{|c|[Format]}{Test2[Inhalt]}
\renewcommand{\arraystretch}{1.3} % Tabellenlinien nicht zu dicht
\usepackage{colortbl}
\arrayrulecolor{gray}	% heller Tabellenlinien

% Nützliches
\usepackage{verbatim}	% u.a. zum auskommentieren via \begin{comment} \end{comment}
\usepackage{tabto}	% Tabs: /tab zum nächsten Tab oder /tabto{.5 \CurrentLineWidth} zur Stelle in der Linie
\NumTabs{6}	% Anzahl von Tabs pro Zeile zum springen
\usepackage{listings} % Source-Code mit Tabs
\usepackage{lstautogobble} 
\usepackage{enumitem}	% Anpassung der enumerates
\setlist[enumerate,1]{label=\arabic*.)}	% global andere Enum-Items
\usepackage{letltxmacro} % neue Definiton von Grundbefehlen
% Nutzen:
%\LetLtxMacro{\oldemph}{\emph}
%\renewcommand{\emph}[1]{\oldemph{#1}}

% Einrichtung von lst
\lstset{
basicstyle=\ttfamily, 
mathescape=true, 
escapeinside=||, 
autogobble, 
tabsize=2,
basicstyle=\footnotesize\sffamily\color{black},
frame=single,
rulecolor=\color{lightgray},
numbers=left,
numbersep=5pt,
%basicstyle=\footnotesize\sffamily\color{black},
commentstyle=\color{gray},
%frame=single,
%numbers=left,
%numbersep=5pt,
numberstyle=\tiny\color{gray},
keywordstyle=\color{green},
%showspaces=false,
showstringspaces=false,
stringstyle=\color{orange},
tabsize=2,
literate=%
    {Ö}{{\"O}}1
    {Ä}{{\"A}}1
    {Ü}{{\"U}}1
    {ß}{{\ss}}1
    {ü}{{\"u}}1
    {ä}{{\"a}}1
    {ö}{{\"o}}1
    {~}{{\textasciitilde}}1
}

% BibTeX
\usepackage[backend=bibtex, bibencoding=ascii]{biblatex}	% BibTeX
\usepackage{makeidx}
%\makeglossary
%\makeindex

% Grafiken
\usepackage{graphicx}
\usepackage{epstopdf}	% eps-Vektorgrafiken einfügen

% pdf-Setup
\usepackage{pdfpages}
\usepackage[bookmarks,%
bookmarksopen=false,% Klappt die Bookmarks in Acrobat aus
colorlinks=true,%
linkcolor=black,%
citecolor=red,%
urlcolor=green,%
]{hyperref}

% Titel, Autor usw. werden vor dem Anfang des Dokuments in einem Rutsch definiert…
\newcommand{\DTitel}[1]{\newcommand{\Dokumententitel}{#1}}
\newcommand{\DUntertitel}[1]{\newcommand{\Dokumentenuntertitel}{#1}}
\newcommand{\DAutor}[1]{\newcommand{\Dokumentenautor}{#1}}
\newcommand{\DNotiz}[1]{\newcommand{\Dokumentennotiz}{#1}}
% … Deswegen folgendes erst Nach Dokumentenbeginn ausführen:
\AtBeginDocument{
	\hypersetup{
		pdfauthor={\Dokumentenautor},
		pdftitle={HTW Dresden | \Dokumententitel - \Dokumentenuntertitel},
	}
	\automark[section]{section}
	\automark*[subsection]{subsection}
	\pagestyle{scrheadings}
	\ihead{\includegraphics[height=1.7em]{../LaTeX_master/HTW-Logo.eps}}
	\ohead{\Dokumententitel}
	\ofoot{\footnotesize{\textcolor{darkgray}{Mitschrift von\\ \Dokumentenautor}}}
	% Titelseite
	\title{\includegraphics[width=0.35\textwidth]{../LaTeX_master/HTW-Logo.eps}\\\vspace{0.5em}
	\Huge\textbf{\Dokumententitel} \\\vspace*{0,5cm}
	\Large \Dokumentenuntertitel \\\vspace*{4cm}}
	\author{\textcolor{darkgray}{Mitschrift von \Dokumentenautor} \vspace*{1cm}\\\Dokumentennotiz}
}

%% EIGENE BEFEHLE

%Farbdefinitionen
\definecolor{red}{RGB}{180,0,0}
\definecolor{green}{RGB}{75,160,0}
\definecolor{blue}{RGB}{0,75,200}
\definecolor{orange}{RGB}{255,128,0}
\definecolor{yellow}{RGB}{255,245,0}
\definecolor{purple}{RGB}{75,0,160}
\definecolor{cyan}{RGB}{0,160,160}
\definecolor{brown}{RGB}{120,60,10}

% Textfarbe ändern
\newcommand{\tred}[1]{\textcolor{red}{#1}}
\newcommand{\tgreen}[1]{\textcolor{green}{#1}}
\newcommand{\tblue}[1]{\textcolor{blue}{#1}}
\newcommand{\torange}[1]{\textcolor{orange}{#1}}
\newcommand{\tyellow}[1]{\textcolor{yellow}{#1}}
\newcommand{\tpurple}[1]{\textcolor{purple}{#1}}
\newcommand{\tcyan}[1]{\textcolor{cyan}{#1}}
\newcommand{\tbrown}[1]{\textcolor{brown}{#1}}

% Umstellen der Tabellen Definition
\newcommand{\mpb}[1][.3]{\begin{minipage}{#1\textwidth}\vspace*{3pt}}
\newcommand{\mpe}{\vspace*{3pt}\end{minipage}}

\newcommand{\resultul}[1]{\underline{\underline{#1}}}
\newcommand{\parskp}{$ $\\}	% new line after paragraph
\newcommand{\corr}{\;\widehat{=}\;}
\newcommand{\mdeg}{^{\circ}}

\newcommand{\nok}[2]{\begin{pmatrix}#1\\#2\end{pmatrix}}	% n über k

\bibliography{../Literatur/HTW_Literatur.bib}

% Definition von Titel, Autor usw.
\DTitel{Technische Informatik}
\DUntertitel{Vorlesungsskript}
\DAutor{Falk-Jonatan Strube}
\DNotiz{Vorlesung von Dr. Boris Hollas}

\begin{document}

\maketitle
\newpage
\tableofcontents
\newpage

\section*{Inhalte}
Grundlage: Grundkurs Theoretische Informatik \cite{hollas2015grundkurs}
\begin{itemize}
\item Formale Sprachen
\begin{itemize}
\item Reguläre Sprachen
\begin{itemize}
\item Endliche Automaten
\item Reguläre Ausdrücke
\end{itemize}
\item Nichtreguläre Sprachen
\item Kontextfreie Sprachen
\begin{itemize}
\item Kellerautomaten
\item Grammatiken
\end{itemize}
\end{itemize}
\item Berechenbarkeit
\begin{itemize}
\item Halteproblem
\end{itemize}
\item Komplexitätsklassen
\begin{itemize}
\item $P$
\item $NP$
\item $NP$-vollständige Probleme
\end{itemize}
\end{itemize}

\section{Automaten und Formale Sprachen}
\subsection{Formale Sprachen}
\paragraph{Def.:} Ein Alphabet ist eine Menge $\Sigma \not = \emptyset$ (Symbole in $\Sigma$ -- müssen nicht einzelne Buchstaben sein, auch Wörter usw. [bspw. „if“ oder „else“ im Alphabet der Programmiersprache C]).

\paragraph{Def.:} Für $w_1, ..., w_n \in \Sigma$ ist $w=w_1...w_n$ ein Wort der Länge $n$.\\
Das Wort $\varepsilon$ ist das \emph{leere Wort}.\\
Die Menge aller Wörter bezeichnen wir mit $\Sigma^*$ (einschließlich dem leeren Wort).

\subparagraph{Bsp.:} $\Sigma = \{a,b,c\}\quad \rightarrow \Sigma^*=\{\varepsilon, a, b, c, aa, ab,a c, aaa, ...\}$
\paragraph{Def.:} Für Wörter $a,b \in \Sigma^*$ ist $ab$ die Konkatenation dieser Wörter.\\
Für ein Wort $w$ ist $w^n$ die $n$-fache Konkatenation von $w$, wobei $w^0=\varepsilon$.

\subparagraph{Bemerkung:} Für alle $w \in \Sigma^*$ gilt $\varepsilon w = w = w \varepsilon$. $\varepsilon$ ist also das neutrale Element der Konkatenation.

\paragraph{Def.:} Eine \emph{formale Sprache} ist eine Teilmenge von $\sigma^*$.

\paragraph{Def.:} Für Sprachen $A, B$ ist $AB=\{ab \;|\; a \in A, \; b \in B\}$ sowie $A^n=\prod_{i=1}^{n}A$, wobei $A^0=\{\varepsilon\}$.
\subparagraph{Bemerkung:} $\emptyset, \varepsilon, \{\varepsilon\}$ sind unterschiedliche Dinge (leere Menge, leeres Wort, Menge mit leerem Wort).

\subparagraph{Bemerkung:} $\Sigma^*$ lässt sich ebenfalls definieren durch $\Sigma^*=\bigcup_{n\geq 0}\Sigma^n$.\\
Ferner ist $\Sigma^+=\Sigma^*-\{\varepsilon\}$.

\subsection{Reguläre Sprachen}
\subsubsection{Deterministische endliche Automaten (DFA)}
\subparagraph{Bsp.:} \parskp
ABB1\\
(Pfeil zeigt auf Startzustand, Endzustand ist doppelt umrandet)\\
Dieser DFA akzeptiert alle Wörter über $\Sigma = \{a,b,c\}$, die $abc$ enthalten.\\
Deterministisch: Es gibt genau ein Folgezustand. Von jedem Knoten aus gibt es genau eine Kante für jedes Zeichen, nicht mehrere und nicht keine.
\subparagraph{Bsp.:} \parskp
ABB2\\
Dieser DFA erkennt die Sprache $\{a,b,c\}^*$.
\paragraph{Def.:} Ein DFA ist ein Tupel $\mathcal{M}=(Z, \Sigma, \delta, z_0, )$
\begin{itemize}
\item $Z$: Menge der Zustände
\item $\Sigma$: Eingabealphabet
\item $\delta$: Überführungsfunktion $Z\times \Sigma \rightarrow Z$. Dabei bedeutet $\delta (z,a) = z'$, dass $\mathcal{M}$ im Zustand $z$ für das Zeichen $a$ in den Zustand $z'$ wechselt.
\item $z_0\in Z$: Startzustand
\item $E$: Menge der Endzustände
\end{itemize}
$\delta$: ABB3

\paragraph{Def.:} Die erweiterte Überführungsfunktion $\hat{\delta}:Z\times \Sigma^*\rightarrow Z$ ist definiert durch\\
$\hat{\delta}(z,w)=\begin{cases}
z & \text{für }w = \varepsilon\\
\hat{\delta}(\delta(z,a),x) & \text{für }w=ax \text{ mit } a\in \Sigma ,x \in \Sigma^*
\end{cases}$\\
Dazu vergleichbarer C-Code:
\begin{lstlisting}[language=C]
int $\hat{\delta}$(int z, char* w){
	if ( strlen(w) == 0 )
		return z;
	else
		return $\hat{\delta}$($\delta$(z, w[0]), w[1]);
\end{lstlisting}
Veranschaulichung:\\
ABB4\\
Die erweiterte Überführungsfunktion bestimmt den Zustand nach dem vollständingen Lesen eines Wortes.

\subparagraph{Bsp.:} \parskp
ABB5
\begin{align*}
\hat{\delta}(z_0, aaba) &= \hat{\delta}(\delta(z_0,a),aba)=\\
\hat{\delta}(z_0, aba) &= \hat{\delta}(\delta(z_0,a),ba)=\\
\hat{\delta}(z_0, ba) &= \hat{\delta}(\delta(z_0,b),a)=\\
\hat{\delta}(z_E, a) &= \hat{\delta}(\delta(z_E,a), \varepsilon)=\\
\hat{\delta}(z_E, \varepsilon) &= z_E\\
\end{align*}
Die von $\mathcal{M}$ \emph{akzeptierte Sprache} ist $L(\mathcal{M})=\{w\in \Sigma^* \;|\; \hat{\delta}(z_0, w) \in E\}$

\newpage
\printbibliography
\end{document}